We shall digress but for a brief moment to introduce functions upon spheres.
A more thorough discussion of the topic immediately at hand is of course to be found in any monograph on the topic.
The purpose of this digression is to clearly illustrate for what reasons and by which means we expand the \textsc{Green} function into representations in terms of spherical functions.
Let $\nu$ be a positive semiïnteger, meaning that $2\nu \in \naturals^{+}$, and let $n$ be a positive integer.
The \textsc{Gegenbauer} ultraspherical polynomials $\gls{gegenbauer}(x)$ then constitute a complete orthonormal basis for the $L^2$-space on the unit interval, with respect to the measure ${(1 - x^2)}^{\nu - \sfrac{1}{2}}\,\dee x$, indexed in $n$.%
\footnote{%
  \cite{takeuchi1994modern} \textsc{Takeuchi}, p.222
}
%
The significance of this measure is not immediately apparent.
Alternatively to the above definition as generating functions, we have the definition through the hypergeometric function,%
\footnote{%
  \cite{takeuchi1994modern} \textsc{Takeuchi}, p.194
}

\begin{equation}
  \gegenbauer_n^{\nu} (x) = \frac{\Gamma(2\nu + n)}{\gls{factorial}\Gamma(2\nu)} \hypergeometric{\Big( 2\nu + n,\, -n;\, \nu + \sfrac{1}{2};\, \sfrac{1}{2} - \sfrac{x}{2} \Big)},
\end{equation}

\noindent
The hypergeometric function $\hypergeometric(a,b;c;\ezh)$ is a special function, whose full treatment is far beyond the scope of this work.
It warrants mentioning that it may either be defined in terms of a power series,%
\footnote{%
  \cite{abramowitz1965handbook} \textsc{Oberhettinger}, eq.15.1.1, p.556
}
or more relevant for our purposes, as the unique real analytic of the hypergeometric equation

\begin{subequations}\label{eq:hypergeometric}
  \begin{equation}\tag{\theequation{a, b}}
    \ezh(1 - \ezh)\frac{\dee^2 f}{\dee\ezh^2} + \big( c - (a+b+1)\ezh \big) \frac{\dee f}{\dee\ezh} - ab f = 0, \qquad \absl{\ezh} < 1.
  \end{equation}
\end{subequations}

\noindent
This equation will reappear in an impending subsection, as a direct consequence of the theory for which we seek a foundation.
The ultraspherical functions are related to the associated \textsc{Legendre} spherical function through the following expression.%
\footnote{%
  \cite{abramowitz1965handbook} \textsc{Hochstrasser}, eq.22.5.60, p.780
}

\begin{equation}
  \gegenbauer_n^{\nu} (x) = \frac{\Gamma(\nu + \sfrac{1}{2}) \Gamma(2\nu + n)}{n! \Gamma(2\nu)} {\left( \frac{x^2 - 1}{4} \right)}^{\frac{1 - 2\nu}{4}} \legendre_{n + \nu - \sfrac{1}{2}}^{\sfrac{1}{2} - \nu} (x).
\end{equation}

\noindent
Through algebraic manipulations, the definition of the associated \textsc{Legendre} functions are found to be%
\footnote{%
  \cite{abramowitz1965handbook} \textsc{Stegun}, eq.8.1.2, p.332
}

\begin{equation}
  \legendre_{n}^{m} (\ezh) = \frac{1}{\Gamma(1 - m)} {\Bigg( \frac{\ezh + 1}{\ezh - 1} \Bigg)}^{\sfrac{m}{2}} \hypergeometric{\Big( -n,\, n + 1;\, 1 - m;\, \sfrac{1}{2} - \sfrac{\ezh}{2} \Big)}
\end{equation}

\noindent
For $\nu = \sfrac{1}{2}$, we have that $m = 0$, where $\legendre_{k}^{0} = \legendre_{k}$ is the $k$\textsuperscript{th} \textsc{Legendre} polynomial.%
\footnote{%
  \cite{abramowitz1965handbook} \textsc{Stegun}, pp.332--353
}
